\section*{Chapter \ref{ch:g9:landforms} --- Landforms}

\begin{solution}{exo:g9:landforms:1}
(a) The wind. (b) Water dissolving the \omterm{def:g3:stones-pebbles-sand:rock}{rock} (and depositing calcite in the
\omterm{def:g9:landforms:karst}{cave}). (c) The \omterm{def:g2:water-on-earth:ocean}{sea} (\omterm{def:g9:landforms:beach}{longshore drift}). (d) Ice. (e) The \omterm{def:g2:water-on-earth:ocean}{sea}. (f) A \omterm{def:g1:landscapes-around-us:water}{river}.
\end{solution}

\begin{solution}{exo:g9:landforms:2}
(a) \omterm{def:g9:landforms:uplift}{Uplift}. (b) The \omterm{def:g9:landforms:dune}{slip face}. (c) A \omterm{def:g9:landforms:karst}{sinkhole}. (d) \omterm{def:g9:landforms:beach}{Longshore drift}.
(e) A \omterm{def:g9:landforms:karst}{karst}.
\end{solution}

\begin{solution}{exo:g9:landforms:3}
(a) False: the horns point downwind; the round nose faces the wind.
(b) True. (c) False: the wind moves it grain by grain, up the gentle slope
and down the \omterm{def:g9:landforms:dune}{slip face}. (d) True.
\end{solution}

\begin{solution}{exo:g9:landforms:4}
(b), (d), (e), (a), (c).
\end{solution}

\begin{solution}{exo:g9:landforms:5}
(a) From the west: the horns point downwind. (b) On the east side, between
the horns. (c) $60 \div 6 = 10$ years. (d) A larger \omterm{def:g9:landforms:dune}{dune} holds more \omterm{def:g3:stones-pebbles-sand:sand}{sand},
and all of it has to be carried over the crest for the \omterm{def:g9:landforms:dune}{dune} to move
forward.
\end{solution}

\begin{solution}{exo:g9:landforms:6}
(a) $7 \times 1 = \qty{7}{m}$. (b) $2 \div 12 = \frac{1}{6}$ of a year.
(c) Seven years of retreat in $\frac{1}{6}$ of a year: $7 \times 6 = 42$
times faster. (d) The waves first cut a notch; the \omterm{def:g3:stones-pebbles-sand:rock}{rock} above hangs over it
for a while, without moving; then the overhang falls all at once, and the
edge jumps back.
\end{solution}

\begin{solution}{exo:g9:landforms:7}
(a) Towards the east. (b) Three strokes up the \omterm{def:g9:landforms:beach}{beach} slanting to the
north-east, each followed by a stroke straight back down to the south: a
zigzag moving east. (c) A \omterm{def:g9:landforms:beach}{spit}, growing eastward across the mouth of the
bay. (d) The \omterm{def:g9:landforms:beach}{spit} shelters the water behind it from the waves; in calm
water, the fine \omterm{def:g3:stones-pebbles-sand:sand}{mud} can settle.
\end{solution}

\begin{solution}{exo:g9:landforms:8}
(a) It takes up carbon dioxide from the air and from the soil. (b) The
weak acid dissolves limestone slowly, but hardly attacks granite. (c) The
water can only enter the \omterm{def:g3:stones-pebbles-sand:rock}{rock} through its cracks and along its layers; it
widens these openings first. (d) The \omterm{def:g3:stones-pebbles-sand:rock}{rock} below the surface is dissolved
and the ground sinks slowly; or the roof of a \omterm{def:g9:landforms:karst}{cave} collapses.
\end{solution}

\begin{solution}{exo:g9:landforms:9}
(a) When the drop meets the air of the \omterm{def:g9:landforms:karst}{cave}, it gives off some of its
carbon dioxide and can no longer hold all of the dissolved limestone, which
comes out as calcite. (b) A stalagmite grows from drops that splash on the
floor and spread out; there is no drop hanging at its tip to build a tube.
(c) $\qty{30}{cm} = \qty{300}{mm}$; $300 \div 0.1 = 3000 = 3 \times 10^3$
years.
\end{solution}

\begin{solution}{exo:g9:landforms:10}
(a) $4600 \div 10^7 = 4.6 \times 10^{-4}$~m a year, or \qty{0.46}{mm} a
year. (b) $0.46 - 0.2 = \qty{0.26}{mm}$ a year. (c) $\qty{1}{m} =
\qty{1000}{mm}$; $1000 \div 0.46 \approx 2200$, about two thousand times
faster. (d) It is kept up for millions of years: $\qty{0.26}{mm} \times
10^7 = 2.6 \times 10^6$~mm, about \qty{2.6}{km}.
\end{solution}

\begin{solution}{exo:g9:landforms:11}
(a) Shape: crescent \omterm{def:g9:landforms:dune}{dunes} with their \omterm{def:g9:landforms:dune}{slip face} on the south-east side.
Agent and process: the wind, blowing from the north-west, \omterm{def:g5:water-shapes-the-land:deposition}{deposition}.
Material: loose \omterm{def:g3:stones-pebbles-sand:sand}{sand}. (b) Shape: a \omterm{def:g9:landforms:cliff}{cliff} above a flat shelf uncovered at low
tide. Agent and process: the \omterm{def:g2:water-on-earth:ocean}{sea}, \omterm{def:g5:water-shapes-the-land:erosion}{erosion}; the shelf is a wave-cut
platform left by the retreating \omterm{def:g9:landforms:cliff}{cliff}. Material: white \omterm{def:g3:stones-pebbles-sand:rock}{rock}, probably
chalk, a soft \omterm{def:g3:stones-pebbles-sand:rock}{rock}. (c) Shape: a stream that vanishes into a hole, a \omterm{def:g6:water-underground:spring}{spring}
far away. Agent and process: water dissolving the \omterm{def:g3:stones-pebbles-sand:rock}{rock} and flowing
underground. Material: limestone; the plateau is a \omterm{def:g9:landforms:karst}{karst}.
\end{solution}

\begin{solution}{exo:g9:landforms:12}
(a) The small \omterm{def:g9:landforms:dune}{dune}: $45 \div 6 = 7.5$ years; the larger one: $45 \div 3 =
15$ years. (b) $45 \div 100 = 0.45$ of a year, about five or six months.
(c) For example, plant grasses on the \omterm{def:g9:landforms:dune}{dunes}: their roots hold the \omterm{def:g3:stones-pebbles-sand:sand}{sand}, so
the wind can no longer take it.
\end{solution}

\begin{solution}{pb:g9:landforms:1}
\textbf{1.} $46.0 - 30.4 = \qty{15.6}{m}$.

\textbf{2.} $15.6 \div 20 = \qty{0.78}{m}$ a year.

\textbf{3.} Between 2012 and 2014: $41.0 - 36.2 = \qty{4.8}{m}$ in two
years, or $4.8 \div 2 = \qty{2.4}{m}$ a year.

\textbf{4.} $2.4 \div 0.78 \approx 3$ times faster. Probably a stormy
winter, during which large overhangs fell.

\textbf{5.} \qty{78}{cm} a year; $0.78 \times 100 = \qty{78}{m}$ per
century.

\textbf{6.} Yes: \qty{0.78}{m} a year is a little less than a metre a year.

\textbf{7.} $17 \div 0.78 \approx 22$ years.

\textbf{8.} The \omterm{def:g9:landforms:cliff}{cliff} retreats in jumps: a two-year interval may hold a
large fall, or none at all. Over 20 years, the jumps and the calm years
average out.

\textbf{9.} The edge has $30.4 - 10 = \qty{20.4}{m}$ to go: $20.4 \div 0.78
\approx 26$ years, so in 2050.

\textbf{10.} $2 \times 0.78 = \qty{1.56}{m}$ a year; $30.4 \div 1.56
\approx 19.5$ years after 2024, so in 2043.

\textbf{11.} The edge of a \omterm{def:g9:landforms:cliff}{cliff} can give way at any moment; standing on
it, or near it, is dangerous.

\textbf{12.} $30.4 \div 0.78 \approx 39$ years after 2024: in
\textbf{2063}.
\end{solution}
